% Lab 1 of Intelligent Robotics (ASY2): kinematics and odometry of a differential and a mecanum drive.
% Build: make  (latexmk -pdf). Figures are TikZ files in figures/, no external graphics.
\documentclass[11pt,a4paper,parskip=half]{scrartcl}
\usepackage[english]{babel}
\usepackage[T1]{fontenc}
\usepackage[utf8]{inputenc}
\usepackage[sfdefault,lining]{FiraSans}
\usepackage[scale=0.88]{FiraMono}
\usepackage[a4paper,left=2.4cm,right=2.4cm,top=2.6cm,bottom=2.6cm,headsep=0.8cm]{geometry}
\usepackage{amsmath,amssymb}
\usepackage{booktabs,array}
\usepackage{xcolor}
\usepackage{listings}
\usepackage{tikz}
\usetikzlibrary{arrows.meta,calc,shapes.geometric}
\usepackage[skins,breakable]{tcolorbox}
\usepackage{fontawesome5}
\usepackage[font=small,labelfont={bf,color=ohmred},skip=4pt]{caption}
\usepackage{enumitem}
\usepackage{microtype}
\usepackage[headsepline=0.4pt]{scrlayer-scrpage}
\usepackage[colorlinks=true,linkcolor=ohmred,urlcolor=ohmred]{hyperref}
\usepackage{cleveref}

% ---------------------------------------------------------------- colours and headings
\definecolor{ohmred}{HTML}{C8372D}
\definecolor{ohmgreen}{HTML}{2E8B57}
\definecolor{ohmamber}{HTML}{D98E04}
\definecolor{ohmgrey}{HTML}{5A5A5A}
\addtokomafont{disposition}{\color{ohmred}}
\addtokomafont{section}{\Large}
\addtokomafont{subsection}{\large}
\setkomafont{pageheadfoot}{\small\color{ohmgrey}}
\setkomafont{pagenumber}{\small\bfseries\color{ohmred}}
\clearpairofpagestyles
\ihead{Intelligent Robotics (ASY2) · Lab 1}
\ohead{Kinematics and odometry}
\ofoot{\pagemark}
\pagestyle{scrheadings}

\lstset{basicstyle=\ttfamily\small, columns=fullflexible, keepspaces=true,
        showstringspaces=false, breaklines=true, frame=none, backgroundcolor=\color{black!5},
        xleftmargin=8pt, xrightmargin=8pt, framexleftmargin=6pt, framexrightmargin=6pt,
        aboveskip=4pt, belowskip=4pt}

\newcommand{\thema}[1]{\texttt{#1}}
\newcommand{\datei}[1]{\texttt{#1}}
\newcommand{\minutes}[1]{\hfill{\normalsize\mdseries\color{ohmgrey}\faClock[regular]~#1 min}}

% ---------------------------------------------------------------- boxes
\tcbset{base/.style={enhanced, breakable, boxrule=0pt, arc=1.5pt, left=8pt, right=8pt,
        top=5pt, bottom=5pt, before skip=8pt, after skip=8pt}}
\newtcolorbox{todo}[1]{base, colback=ohmred!5, borderline west={2.5pt}{0pt}{ohmred},
        title={\faPen~#1}, coltitle=ohmred, colbacktitle=ohmred!5, fonttitle=\bfseries,
        attach title to upper={\quad}}
\newtcolorbox{script}[1]{base, colback=black!3, borderline west={2.5pt}{0pt}{ohmgrey},
        title={\faBook~Lecture script, #1}, coltitle=ohmgrey, colbacktitle=black!3,
        fonttitle=\small\bfseries, toptitle=2pt, bottomtitle=0pt}
\newtcolorbox{note}{base, colback=ohmamber!8, borderline west={2.5pt}{0pt}{ohmamber}}
\newtcolorbox{checkpoint}{base, colback=ohmgreen!7, borderline west={2.5pt}{0pt}{ohmgreen},
        title={\faCheckCircle[regular]~Checkpoint}, coltitle=ohmgreen, colbacktitle=ohmgreen!7,
        fonttitle=\bfseries, attach title to upper={\quad}}

\setlist[enumerate,1]{leftmargin=*,itemsep=3pt,topsep=2pt}
\setlist[itemize,1]{leftmargin=*,itemsep=2pt,topsep=2pt}

\begin{document}
\sloppy

% ---------------------------------------------------------------- title
\begin{tcolorbox}[enhanced, colback=ohmred, colframe=ohmred, arc=0pt, left=14pt, right=14pt,
                  top=12pt, bottom=12pt, fontupper=\color{white}]
  {\small Intelligent Robotics (ASY2) · Lab 1}\\[6pt]
  {\LARGE\bfseries Kinematics and odometry}\\[4pt]
  {\large A differential drive and a mecanum drive as ROS\,2 nodes}
\end{tcolorbox}

\noindent
\begin{tabular}{@{}>{\color{ohmgrey}}l l@{\hspace{2.4em}}>{\color{ohmgrey}}l l@{}}
  \faClock[regular]~Duration & 180 minutes & \faUsers~Work in & pairs \\
  \faBook~Script & Cycle 1, sections 1.1 and 1.2 & \faFolderOpen[regular]~Files & \datei{ohm-odometry/student/}
\end{tabular}

\section{Overview}
You write two ROS\,2 nodes for a simulated robot, one for a differential drive and one for a
mecanum drive. Each node does three things:
\begin{itemize}
  \item it turns the keyboard command into wheel speeds (inverse kinematics),
  \item it turns the measured wheel speeds back into the speed of the robot (forward kinematics),
  \item it adds up these speeds to a position (odometry).
\end{itemize}
The node talks to the simulator only through ROS topics. Which topics these are, you find out
yourself.

\begin{table}[ht]\centering
\caption{Schedule.}
\label{tab:schedule}
\begin{tabular}{@{}rll@{}}
\toprule
min & task & done when \\
\midrule
0--20 & \cref{sec:task0} Explore the simulator & \cref{tab:interface} is filled in \\
20--55 & \cref{sec:task1} Differential drive: kinematics & the robot drives with the keys \\
55--85 & \cref{sec:task2} Differential drive: odometry & the green outline follows the robot \\
85--95 & break & \\
95--130 & \cref{sec:task3} Mecanum drive: kinematics & the robot drives sideways with the keys \\
130--155 & \cref{sec:task4} Mecanum drive: odometry & the green outline follows the robot \\
155--180 & \cref{sec:extra} Extra tasks & \\
\bottomrule
\end{tabular}
\end{table}

\paragraph{Notation} as in the script: pose \(\mathbf p = (x, y, \phi)^{\mathsf T}\), body twist
\(\xi = (v_x, v_y, \omega)^{\mathsf T}\). The \(x\)-axis points forward, \(y\) to the \textbf{left}, and the
heading \(\phi\) counts counter-clockwise. Speeds are in \(\mathrm{m/s}\) and \(\mathrm{rad/s}\), wheel
speeds in \(\mathrm{rad/s}\). The equations you need are repeated in the tasks, with the numbers
they have in the script.

\newpage
\section{Getting started}
\paragraph{Your own copy.} Fork \datei{ohm-odometry} on GitHub (button \emph{Fork} at the top right),
then clone your fork and the simulator into the \datei{src/} folder of your workspace:
\begin{lstlisting}
git clone git@github.com:<your-name>/ohm-odometry.git
\end{lstlisting}
\begin{lstlisting}
git clone git@github.com:christianpfitzner/ir-mecanum-lab.git
\end{lstlisting}
Your work then lives in your own repository. Commit after every task and push, so nothing is lost at
the end of the session.

\paragraph{Build.} Build both once, then source the workspace in every new terminal:
\begin{lstlisting}
colcon build --symlink-install --packages-select mecanum_lab_interfaces mecanum_lab ohm_odometry
\end{lstlisting}
\begin{lstlisting}
source install/setup.bash
\end{lstlisting}
You do not need to build again after a change to your node. Just start it again.

\paragraph{Two ways to start the simulator.}
\begin{itemize}
  \item \thema{diff\_demo.launch.py} and \thema{mecanum\_demo.launch.py}: the simulator uses its own
        kinematics. The keys drive the robot, no node of yours is needed.
  \item \thema{diff.launch.py} and \thema{mecanum.launch.py}: the robot moves only on the wheel speeds
        your node sends. Without your node it stands still, with wrong kinematics it drives wrong.
\end{itemize}

\paragraph{Your files.} \datei{student/diff\_drive\_node.py} for tasks 1 and 2,
\datei{student/mecanum\_drive\_node.py} for tasks 3 and 4. Each is a complete ROS\,2 node. You only
change the lines marked \thema{TODO}; subscribing, publishing and reading the messages are already
there. Your node runs in a second terminal:
\begin{lstlisting}
ros2 run ohm_odometry diff_drive_node.py
\end{lstlisting}

\paragraph{What the node does.} 50 times per second the simulator sends the wheel encoders. For each
message the node computes the wheel speeds, calls your forward kinematics and your odometry step,
and publishes the new pose. The simulator draws it as a green outline of the robot, with a line for
the heading and a green line for the path. Also 50 times per second, the node takes the last keyboard
command, calls your inverse kinematics and publishes the wheel speeds.

\begin{figure}[ht]
  \centering
  % TikZ figure: the two nodes and the five topics between them, drawn like rqt_graph
% (nodes as ellipses, topics as boxes). Topic names are left out on purpose: they are task 0.
\begin{tikzpicture}[
  font=\small,
  node/.style={ellipse, draw=black!70, line width=0.8pt, fill=ohmred!7, minimum height=1.3cm,
               minimum width=3.0cm, align=center},
  topic/.style={rectangle, draw=black!55, fill=black!3, minimum width=4.6cm, minimum height=0.95cm,
                align=center, inner sep=2pt},
  flow/.style={-{Stealth[length=2.2mm]}, line width=0.7pt, draw=black!70},
]
\node[node] (sim) at (0, 0) {\texttt{/mecanum\_sim}\\[-1pt]{\footnotesize simulator + window}};
\node[node, fill=ohmgreen!12] (me) at (11.0, 0) {\texttt{/diff\_drive}\\[-1pt]{\footnotesize your node}};

\node[topic] (enc) at (5.5,  2.6) {wheel encoders\\[-1pt]{\scriptsize\texttt{sensor\_msgs/JointState}, 50 Hz}};
\node[topic] (key) at (5.5,  1.4) {keyboard command\\[-1pt]{\scriptsize\texttt{geometry\_msgs/Twist}}};
\node[topic] (tru) at (5.5,  0.2) {true pose\\[-1pt]{\scriptsize\texttt{geometry\_msgs/PoseStamped}}};
\node[topic] (whl) at (5.5, -1.2) {wheel speeds\\[-1pt]{\scriptsize\texttt{std\_msgs/Float64MultiArray}, 50 Hz}};
\node[topic] (est) at (5.5, -2.4) {your pose (green outline)\\[-1pt]
  {\scriptsize\texttt{geometry\_msgs/PoseWithCovarianceStamped}}};

\foreach \t in {enc, key, tru} {
  \draw[flow] (sim) -- (\t.west);
  \draw[flow] (\t.east) -- (me);
}
\foreach \t in {whl, est} {
  \draw[flow] (me) -- (\t.east);
  \draw[flow] (\t.west) -- (sim);
}
\node[font=\footnotesize, text=black!60, align=center] at (5.5, 3.55) {topic names: see task 0};
\end{tikzpicture}

  \caption{The simulator and your node, drawn like \thema{rqt\_graph}: nodes are ellipses, topics are
    boxes, arrows point from publisher to subscriber.}
  \label{fig:graph}
\end{figure}

\newpage
\section{Task 0: explore the simulator\minutes{20}}
\label{sec:task0}
\begin{enumerate}
  \item Start the simulator with its own kinematics:
\begin{lstlisting}
ros2 launch ohm_odometry diff_demo.launch.py
\end{lstlisting}
        Click into the window and drive with \thema{w}/\thema{s} (forward, back) and
        \thema{q}/\thema{e} (turn). What happens with \thema{a}/\thema{d} (sideways)?
  \item In a second terminal, list the topics and look at each one in \cref{tab:interface}:
\begin{lstlisting}
ros2 topic list
\end{lstlisting}
\begin{lstlisting}
ros2 topic info -v <topic>
\end{lstlisting}
  \item Drive forward and watch the encoder message. Which numbers change? How often does it arrive?
\begin{lstlisting}
ros2 topic echo <topic>
\end{lstlisting}
\begin{lstlisting}
ros2 topic hz <topic>
\end{lstlisting}
  \item Open \thema{rqt\_graph} and switch off the filters \emph{Dead sinks} and \emph{Leaf topics}.
        Compare with \cref{fig:graph}. What is still missing?
  \item Stop the simulator (Ctrl-C) and start \thema{mecanum\_demo.launch.py}. What changes in the
        encoder message?
\end{enumerate}

\begin{table}[ht]\centering
\renewcommand{\arraystretch}{2.0}
\caption{The topics between your node and the simulator. Fill in.}
\label{tab:interface}
\begin{tabular}{@{}>{\raggedright\arraybackslash}p{4.2cm}p{3.9cm}p{4.4cm}p{1.3cm}@{}}
\toprule
& topic name & message type & rate \\
\midrule
wheel encoders & & & \\
keyboard command & & & \\
true pose of the robot & & & \\
wheel speeds (your node) & & & \\
pose of your odometry & & & \\
\bottomrule
\end{tabular}
\end{table}

\begin{note}
The last two topics have no publisher yet, the simulator only listens to them. They still show up in
\thema{ros2 topic list}. Your pose goes to the topic of type \thema{PoseWithCovarianceStamped}.
\end{note}

\newpage
\section{Task 1: differential drive, kinematics\minutes{35}}
\label{sec:task1}
\begin{figure}[ht]
  \centering
  % TikZ figure: differential drive from above, inside the world frame.
% Signs as in CONTRACT 5: x forward, y left, phi counter-clockwise; left wheel at y = +b/2.
% Everything of the robot is a path, not a node: paths turn with the scope, nodes would not.
\begin{tikzpicture}[
  scale=0.95, every node/.style={font=\small},
  axis/.style={-{Stealth[length=2.6mm]}, line width=1pt},
  vel/.style={-{Stealth[length=2mm]}, draw=ohmred, line width=0.9pt},
  dim/.style={{Stealth[length=1.5mm]}-{Stealth[length=1.5mm]}, draw=black!55,
              line width=0.5pt, font=\footnotesize},
]
% world frame
\draw[axis, black!60] (-0.4,0) -- (9.4,0) node[below left] {$x_W$};
\draw[axis, black!60] (0,-0.4) -- (0,6.0) node[below left] {$y_W$};
\node[below left, font=\footnotesize, text=black!60] at (0,0) {$W$};

% pose of the robot, drawn first so the robot covers the dashed lines
\coordinate (P) at (4.4,2.6);
\draw[dashed, black!40] (P) -- (P |- 0,0) node[below, text=black] {$x$};
\draw[dashed, black!40] (P) -- (0,0 |- P) node[left, text=black] {$y$};
\draw[dashed, black!40] ($(P)+(1.7,0)$) -- ++(2.0,0);       % heading reference, clear of the robot
\draw[-{Stealth[length=2mm]}, black!70] ($(P)+(3.2,0)$) arc (0:30:3.2)
  node[pos=0.5, right=1pt] {$\phi$};

\begin{scope}[shift={(P)}, rotate=30]
  \draw[line width=0.8pt, rounded corners=3pt, fill=black!4] (-0.9,-1.0) rectangle (1.2,1.0);
  \draw[line width=0.7pt, fill=black!30] (-0.55, 1.08) rectangle (0.55, 1.42);   % left wheel
  \draw[line width=0.7pt, fill=black!30] (-0.55,-1.42) rectangle (0.55,-1.08);   % right wheel
  \draw[black!60] (0.95,0) circle (0.12);                                        % caster
  \fill[black] (0,0) circle (1.4pt);
  \draw[axis] (0,0) -- (2.4,0) node[right] {$x_B$};
  \draw[axis] (0,0) -- (0,2.4) node[above] {$y_B$};
  \draw[vel] (0.6, 1.25) -- (1.4, 1.25) node[right, font=\footnotesize] {$v_l$};
  \draw[vel] (0.6,-1.25) -- (2.0,-1.25) node[below right, font=\footnotesize] {$v_r$};
  \node[font=\footnotesize] at (-0.8, 1.25) {L};
  \node[font=\footnotesize] at (-0.8,-1.25) {R};
  \draw[dim] (-1.3,-1.25) -- (-1.3,1.25) node[midway, fill=white, inner sep=1pt] {$b$};
\end{scope}
\end{tikzpicture}

  \caption{Differential drive at the pose \((x, y, \phi)\). The wheels sit \(b/2\) left and right of the
    centre. Compare figure 1.4 of the script.}
  \label{fig:diff}
\end{figure}

Stop the demo and start the simulator for your own node. The robot no longer reacts to the keys.
\begin{lstlisting}
ros2 launch ohm_odometry diff.launch.py
\end{lstlisting}

\begin{script}{section 1.2.2}
With the speeds \(v_l, v_r\) at which the left and the right wheel roll, and the track \(b\):
\begin{equation}
  v = \frac{v_r + v_l}{2}, \qquad \omega = \frac{v_r - v_l}{b}
  \qquad\Rightarrow\qquad
  v_l = v - \omega\,\frac b2, \qquad v_r = v + \omega\,\frac b2 .
  \tag{1.30}
\end{equation}
\end{script}

\begin{todo}{TODO 1}
Write the topic names of \cref{tab:interface} into the top of \datei{student/diff\_drive\_node.py}.
\end{todo}

\begin{todo}{TODO 2: inverse kinematics}
\thema{inverse\_kinematics(v, omega)} with (1.30). It gives \(v_l, v_r\) in \(\mathrm{m/s}\), the
simulator wants the wheels to turn at \(w\) in \(\mathrm{rad/s}\). A wheel with radius \(r_w\) that
rolls at \(v\) turns at \(w = v / r_w\). In the code, \(b\) is \thema{B} and \(r_w\) is \thema{R}.
\end{todo}

Start your node and drive. Forward: both wheels turn forward at the same speed. \thema{q} (turn
left): the left wheel turns backward, the right one forward. \thema{ros2 topic echo} on your
wheel-speed topic shows the numbers.

\begin{todo}{TODO 3: forward kinematics}
\thema{forward\_kinematics(w\_left, w\_right)}: the left half of (1.30). Turn the wheel speeds back
into rolling speeds first.
\end{todo}

\begin{checkpoint}
The robot drives forward with \thema{w} and turns left with \thema{q}. \thema{rqt\_graph} looks like
\cref{fig:graph}.
\end{checkpoint}

\newpage
\section{Task 2: differential drive, odometry\minutes{30}}
\label{sec:task2}
Each encoder message is one step of \(\Delta t\) seconds. Your forward kinematics gives \(v\) and
\(\omega\) for it; the odometry adds up the steps to a pose.

\begin{script}{section 1.2.2}
Distance of the wheels and of the kinematic centre:
\begin{equation}
  \Delta s_l = v_l\,\Delta t, \qquad \Delta s_r = v_r\,\Delta t, \qquad
  \Delta s_k = \frac{\Delta s_l + \Delta s_r}{2}.
  \tag{1.27}
\end{equation}
Radius of the arc and change of the heading:
\begin{equation}
  r = \frac b2 \cdot \frac{\Delta s_r + \Delta s_l}{\Delta s_r - \Delta s_l}, \qquad
  \Delta\phi = \frac{\Delta s_r - \Delta s_l}{b}.
  \tag{1.28}
\end{equation}
Pose update, with \(\phi_0\) the heading before the step:
\begin{equation}
  \Delta s = 2r\sin\!\Big(\frac{\Delta\phi}{2}\Big), \qquad
  \Delta x = \Delta s\cos\!\Big(\frac{\Delta\phi}{2} + \phi_0\Big), \qquad
  \Delta y = \Delta s\sin\!\Big(\frac{\Delta\phi}{2} + \phi_0\Big).
  \tag{1.29}
\end{equation}
\end{script}

\begin{todo}{TODO 4: odometry step}
\thema{odometry\_step(x, y, phi, v, omega, dt)} with (1.27) to (1.29). In the function you have \(v\)
and \(\omega\), so \(\Delta s_k = v\,\Delta t\), \(\Delta\phi = \omega\,\Delta t\) and
\(r = \Delta s_k / \Delta\phi\).
\end{todo}

\begin{note}
When the robot drives straight, \(\Delta\phi = 0\) and the radius is infinite: Python divides by zero.
What is \(\Delta s\) in that case? Handle it with an \thema{if}.
\end{note}

Start your node and drive. Press \thema{t} in the window to show the true path of the robot next to
the green path of your odometry. While the robot stands still, compare the two poses:
\begin{lstlisting}
ros2 topic echo <true pose topic> --once
\end{lstlisting}
\begin{lstlisting}
ros2 topic echo <your pose topic> --once
\end{lstlisting}

\begin{checkpoint}
After a minute of driving in curves and straight lines, the green outline still covers the robot.
\end{checkpoint}

\newpage
\section{Task 3: mecanum drive, kinematics\minutes{35}}
\label{sec:task3}
Stop the simulator and your node, then start both for the mecanum robot (the demo is
\thema{mecanum\_demo.launch.py}):
\begin{lstlisting}
ros2 launch ohm_odometry mecanum.launch.py
\end{lstlisting}
\begin{lstlisting}
ros2 run ohm_odometry mecanum_drive_node.py
\end{lstlisting}
TODO 1, the topic names, is the same as in task 1.

\begin{figure}[ht]
  \centering
  \resizebox{0.62\linewidth}{!}{% TikZ figure: mecanum drive from above, inside the world frame — the companion of diffdrive.tikz.
% Signs as in CONTRACT 5: x forward, y left, phi counter-clockwise. Wheel order and roller axes as in
% the simulator (mecanum_lab/render.py, ROLLERS): FL and RR rollers along (1,1), FR and RL along (1,-1).
% Everything of the robot is a path, not a node: paths turn with the scope, nodes would not.
\begin{tikzpicture}[
  scale=0.85, every node/.style={font=\small},
  axis/.style={-{Stealth[length=2.6mm]}, line width=1pt},
  vel/.style={-{Stealth[length=2mm]}, draw=ohmred, line width=0.9pt},
  dim/.style={{Stealth[length=1.5mm]}-{Stealth[length=1.5mm]}, draw=black!55,
              line width=0.5pt, font=\footnotesize},
]
% world frame
\draw[axis, black!60] (-0.4,0) -- (10.6,0) node[below left] {$x_W$};
\draw[axis, black!60] (0,-0.4) -- (0,7.4) node[below left] {$y_W$};
\node[below left, font=\footnotesize, text=black!60] at (0,0) {$W$};

% pose of the robot
\coordinate (P) at (4.8,3.5);
\draw[dashed, black!40] (P) -- (P |- 0,0) node[below, text=black] {$x$};
\draw[dashed, black!40] (P) -- (0,0 |- P) node[left, text=black] {$y$};
\draw[dashed, black!40] ($(P)+(3.0,0)$) -- ++(2.2,0);       % heading reference, clear of the robot
\draw[-{Stealth[length=2mm]}, black!70] ($(P)+(4.6,0)$) arc (0:30:4.6)
  node[pos=0.5, right=1pt] {$\phi$};

% one wheel: centre (#1,#2), roller slope #3 (+1: rollers along (1,1), -1: along (1,-1))
\newcommand{\mwheel}[3]{%
  \draw[line width=0.7pt, fill=black!25] ($(#1,#2)+(-0.45,-0.19)$) rectangle ($(#1,#2)+(0.45,0.19)$);
  \begin{scope}
    \clip ($(#1,#2)+(-0.45,-0.19)$) rectangle ($(#1,#2)+(0.45,0.19)$);
    \foreach \t in {-0.6,-0.4,...,0.6}
      \draw[black!70, line width=0.5pt] ($(#1,#2)+(\t-0.22,-0.22*#3)$) -- ($(#1,#2)+(\t+0.22,0.22*#3)$);
  \end{scope}}

\begin{scope}[shift={(P)}, rotate=30, scale=1.3]
  \draw[line width=0.8pt, rounded corners=3pt, fill=black!4] (-1.65,-0.82) rectangle (1.65,0.82);
  \mwheel{ 1.25}{ 1.03}{ 1}    % FL
  \mwheel{ 1.25}{-1.03}{-1}    % FR
  \mwheel{-1.25}{ 1.03}{-1}    % RL
  \mwheel{-1.25}{-1.03}{ 1}    % RR
  \node[font=\footnotesize] at ( 1.25, 1.62) {FL\,$\omega_0$};
  \node[font=\footnotesize] at ( 1.25,-1.62) {FR\,$\omega_1$};
  \node[font=\footnotesize] at (-1.25, 1.62) {RL\,$\omega_2$};
  \node[font=\footnotesize] at (-1.25,-1.62) {RR\,$\omega_3$};
  \fill[black] (0,0) circle (1.4pt);
  \draw[axis] (0,0) -- (2.7,0) node[right] {$x_B$};
  \draw[axis] (0,0) -- (0,2.3) node[above] {$y_B$};
  \draw[vel] (0.25,-0.4) -- (1.05,-0.4) node[right, font=\footnotesize] {$v_x$};
  \draw[vel] (-0.95,0.05) -- (-0.95,0.6) node[right, font=\footnotesize] {$v_y$};
  \draw[vel] (0.34,0.1) arc (16:164:0.35) node[pos=0.5, above, font=\footnotesize] {$\omega$};
  \draw[dim] (-1.25,-2.05) -- (1.25,-2.05) node[midway, fill=white, inner sep=1pt] {$2\,l_x$};
  \draw[dim] (-2.15,-1.03) -- (-2.15,1.03) node[midway, fill=white, inner sep=1pt] {$2\,l_y$};
\end{scope}
\end{tikzpicture}
}
  \caption{Mecanum drive at the pose \((x, y, \phi)\), rollers in X layout. Compare figure 1.11 of the
    script.}
  \label{fig:mecanum}
\end{figure}

\begin{script}{section 1.2.6}
Wheels \(\omega_0\) front left, \(\omega_1\) front right, \(\omega_2\) rear left, \(\omega_3\) rear right,
wheel radius \(r\), half distances \(l_x, l_y\) (\thema{R}, \thema{LX}, \thema{LY} in the code).
\begin{equation}
  \begin{pmatrix}\omega_0\\ \omega_1\\ \omega_2\\ \omega_3\end{pmatrix} = \frac 1r
  \begin{pmatrix}
     1 & -1 & -(l_x + l_y)\\
    -1 & -1 & -(l_x + l_y)\\
     1 &  1 & -(l_x + l_y)\\
    -1 &  1 & -(l_x + l_y)
  \end{pmatrix}
  \begin{pmatrix} v_x\\ v_y\\ \omega \end{pmatrix}
  \tag{1.39}
\end{equation}
\begin{equation}
  \begin{pmatrix} v_x\\ v_y\\ \omega \end{pmatrix} = \frac r4
  \begin{pmatrix}
     1 & -1 & 1 & -1\\
    -1 & -1 & 1 &  1\\
    \frac{-1}{l_x + l_y} & \frac{-1}{l_x + l_y} & \frac{-1}{l_x + l_y} & \frac{-1}{l_x + l_y}
  \end{pmatrix}
  \begin{pmatrix}\omega_0\\ \omega_1\\ \omega_2\\ \omega_3\end{pmatrix}
  \tag{1.40}
\end{equation}
\end{script}

\newpage
\begin{note}
\textbf{The right wheels count the other way round.} In the script, \(\omega_1\) and \(\omega_3\) are
negative when the robot drives forward, as on a robot whose right motors are mounted mirrored. In the
simulator every wheel is positive when it rolls forward. So
\(w_\text{FR} = -\omega_1\) and \(w_\text{RR} = -\omega_3\).
\end{note}

\begin{todo}{TODO 2: inverse kinematics}
\thema{inverse\_kinematics(vx, vy, omega)} with (1.39) and the note above. The line for the front left
wheel is already there.
\end{todo}

\begin{todo}{TODO 3: forward kinematics}
\thema{forward\_kinematics(w\_fl, w\_fr, w\_rl, w\_rr)} with (1.40). Use \(\omega_1 = -w_\text{FR}\) and
\(\omega_3 = -w_\text{RR}\).
\end{todo}

\begin{checkpoint}
Forward, sideways (\thema{a}/\thema{d}) and turning each work without the other two motions.
\end{checkpoint}

\newpage
\section{Task 4: mecanum drive, odometry\minutes{25}}
\label{sec:task4}
The mecanum robot also moves sideways. Its displacement \((v_x\,\Delta t,\; v_y\,\Delta t)\) is given
in the frame of the robot and has to be rotated into the frame of the hall with \(\mathbf R(\phi)\),
the rotation matrix (1.6) of the script. Taking the heading in the middle of the step,
\(\bar\phi = \phi_k + \Delta\phi/2\) with \(\Delta\phi = \omega\,\Delta t\), as in (1.29):
\begin{equation}
  \begin{pmatrix} x \\ y \end{pmatrix}_{k+1} =
  \begin{pmatrix} x \\ y \end{pmatrix}_{k} +
  \mathbf R(\bar\phi) \begin{pmatrix} v_x \\ v_y \end{pmatrix} \Delta t ,
  \qquad
  \phi_{k+1} = \phi_k + \Delta\phi ,
  \qquad
  \mathbf R(\phi) = \begin{pmatrix}\cos\phi & -\sin\phi\\ \sin\phi & \cos\phi\end{pmatrix}.
  \label{eq:odom}
\end{equation}
For \(v_y = 0\) this is the step of task 2. (The exact step along the arc is shorter by the factor
\(\sin(\Delta\phi/2)/(\Delta\phi/2)\). At 50 steps per second that factor differs from 1 by less than
\(10^{-4}\).)

\begin{todo}{TODO 4: odometry step}
\thema{odometry\_step(x, y, phi, vx, vy, omega, dt)}: multiply out \cref{eq:odom} and write the three
lines for \(x\), \(y\) and \(\phi\).
\end{todo}

\begin{checkpoint}
Drive forward, sideways and in curves. The green outline still covers the robot.
\end{checkpoint}

\newpage
\section{Extra tasks}
\label{sec:extra}
Both tasks drive the robot with your own odometry instead of the keyboard. Each has its own node,
\datei{student/square\_node.py} and \datei{student/goto\_node.py}. They take the topic names, the
kinematics and the odometry of your drive node from tasks 1 to 4 and only replace the keyboard by a controller in \thema{command()};
\thema{self.pose} holds \((x, y, \phi)\) of your odometry. Set \thema{ROBOT} at the top of the file to
\thema{"diff"} or \thema{"mecanum"}. Use \(\operatorname{wrap}(\alpha) = \operatorname{atan2}(\sin\alpha,
\cos\alpha)\) to keep angles in \((-\pi, \pi]\).

Start the matching simulator and the node instead of your drive node, for example:
\begin{lstlisting}
ros2 run ohm_odometry square_node.py
\end{lstlisting}

\subsection*{E1: drive a square by odometry}
TODO E1a and E1b in \datei{square\_node.py}. Side length \(a = 1\,\mathrm{m}\). Store the pose \((x_s, y_s, \phi_s)\) at the start of every side
and every corner.
\begin{itemize}
  \item \textbf{Differential drive.} Side: \(v = 0.2\,\mathrm{m/s}\), \(\omega = 0\), until
        \begin{equation}
          d = \sqrt{(x - x_s)^2 + (y - y_s)^2} \;\ge\; a .
          \label{eq:side}
        \end{equation}
        Corner: \(v = 0\), \(\omega = 0.5\,\mathrm{rad/s}\), until
        \(\left|\operatorname{wrap}(\phi - \phi_s)\right| \ge \pi/2\). Four sides, four corners.
  \item \textbf{Mecanum drive.} The robot needs no corners. Drive the four sides with
        \begin{equation*}
          (v_x, v_y) = (0.2, 0),\quad (0, 0.2),\quad (-0.2, 0),\quad (0, -0.2)\ \mathrm{m/s},
        \end{equation*}
        each until \cref{eq:side} holds.
\end{itemize}
How far from its start does the robot end, measured with the true pose? Where does the error come from?

\subsection*{E2: drive to a point}
TODO E2a to E2c in \datei{goto\_node.py}. Goal \(\mathbf g = (x_g, y_g)\), here 2\,m in front of and
1\,m to the left of the start.
Error in the hall:
\begin{equation}
  \Delta x = x_g - x, \qquad \Delta y = y_g - y, \qquad \rho = \sqrt{\Delta x^2 + \Delta y^2}.
\end{equation}
Stop when \(\rho < 0.05\,\mathrm{m}\).
\begin{itemize}
  \item \textbf{Differential drive.} Angle between the heading and the direction to the goal:
        \begin{equation}
          \alpha = \operatorname{wrap}\!\big(\operatorname{atan2}(\Delta y, \Delta x) - \phi\big), \qquad
          v = k_\rho\,\rho\cos\alpha, \qquad \omega = k_\alpha\,\alpha,
        \end{equation}
        with \(k_\rho = 0.5\,\mathrm{s^{-1}}\), \(k_\alpha = 1.5\,\mathrm{s^{-1}}\) and \(v\) limited to
        \(0.3\,\mathrm{m/s}\). The factor \(\cos\alpha\) slows the robot down while it still points
        away from the goal.
  \item \textbf{Mecanum drive.} Rotate the error into the frame of the robot with
        \(\mathbf R^{-1} = \mathbf R^{\mathsf T}\) (script 1.1.6) and drive straight at the goal:
        \begin{equation}
          \begin{pmatrix} e_x \\ e_y \end{pmatrix} = \mathbf R(\phi)^{\mathsf T}
          \begin{pmatrix} \Delta x \\ \Delta y \end{pmatrix}, \qquad
          v_x = k_\rho\,e_x, \qquad v_y = k_\rho\,e_y, \qquad \omega = 0 .
        \end{equation}
        Limit \(\sqrt{v_x^2 + v_y^2}\) to \(0.3\,\mathrm{m/s}\) by scaling both components.
\end{itemize}
Compare the paths of both robots to the same goal.

\end{document}
